\documentclass[letterpaper]{article}
\usepackage{common/ohpc-doc}
\setcounter{secnumdepth}{5}
\setcounter{tocdepth}{5}

% Include git variables
\input{vc.tex}

% Define Base OS and other local macros
\newcommand{\baseOS}{openEuler 22.03}
\newcommand{\OSRepo}{openEuler\_22.03}
\newcommand{\OSTree}{openEuler\_2203}
\newcommand{\OSTag}{oe2203}
\newcommand{\baseos}{openeuler22.03}
\newcommand{\baseosshort}{openeuler22.03}
\newcommand{\provisioner}{Warewulf}
\newcommand{\provheader}{\provisioner{}}
\newcommand{\rms}{OpenPBS}
\newcommand{\rmsshort}{OpenPBS}
\newcommand{\arch}{aarch64}

% Define package manager commands
\newcommand{\pkgmgr}{yum}
\newcommand{\addrepo}{wget -P /etc/yum.repos.d}
\newcommand{\chrootaddrepo}{wget -P \$CHROOT/etc/yum.repos.d}
\newcommand{\clean}{yum clean expire-cache}
\newcommand{\chrootclean}{yum --installroot=\$CHROOT clean expire-cache}
\newcommand{\install}{yum -y install}
\newcommand{\chrootinstall}{yum -y --installroot=\$CHROOT install}
\newcommand{\groupinstall}{yum -y groupinstall}
\newcommand{\groupchrootinstall}{yum -y --installroot=\$CHROOT groupinstall}
\newcommand{\remove}{yum -y remove}
\newcommand{\upgrade}{yum -y upgrade}
\newcommand{\chrootupgrade}{yum -y --installroot=\$CHROOT upgrade}
\newcommand{\tftppkg}{syslinux-tftpboot}

% boolean for os-specific formatting
\toggletrue{isaarch}
\toggletrue{isCentOS}
\toggletrue{isCentOS_ww_pbs_aarch}
\toggletrue{ispbs}
\toggletrue{isWarewulf}

\begin{document}
\graphicspath{{common/figures/}}
\thispagestyle{empty}

% Title Page --------------------------------------------------------
% Title page and running header definition

\lhead{ \small {\color{logodarkgrey}\fontfamily{phv}\selectfont { Install Guide
    (v\OHPCVersion{})}:  {\baseOS{}/\arch{} + \provheader{} + \rms{}} } \vspace*{0pt} }

{\hspace*{4in} \includegraphics[width=1.7in]{ohpc_logo_blue.pdf}}

\vspace*{2cm}
\noindent {\LARGE \color{logodarkgrey} \fontfamily{phv}\selectfont OpenHPC (v\OHPCVersion{})} \vspace*{0.1cm} \\
\noindent {\LARGE \color{logodarkgrey} \fontfamily{phv}\selectfont Cluster Building Recipes} \\ 

{\color{logoblue}\noindent\rule{6.15in}{1.2pt}} \\ 

\noindent {\Large \color{logodarkgrey} \fontfamily{phv}\selectfont \baseOS{} Base OS} \\ 

\noindent{\Large\color{logodarkgrey}\fontfamily{phv}\selectfont{\provisioner{}/\rms{}
Edition for Linux*} (\arch{})} \\

{\color{logoblue}\noindent\rule{6.15in}{1.2pt}} \\ \vspace{0.2cm}  

\iftoggleverb{isxCATstateful}
\vspace*{-0.2cm}
\noindent{\large\color{logodarkgrey}\fontfamily{phv}\selectfont{Stateful Provisioning}}
\fi

\vspace*{2.5in}

\noindent \includegraphics[width=0.75in]{3x_icon.png}

%\vspace*{.5in}
\noindent{\small \color{black} Document Last Update: \VCDateISO} \vspace*{0.1cm} \\ 
{\small \color{black} Document Revision: \VCRevision} \\ \vspace*{0.1cm}

% Disclaimer
\newpage

\vspace*{3.0cm}
\noindent {\Large \color{logoblue} \fontfamily{phv}\selectfont Legal Notice} \\ 

\vspace*{0.5cm}

\noindent Copyright {\small\copyright} 2016-2022, OpenHPC, a Linux Foundation
Collaborative Project. All rights reserved. \\

\vspace*{0.1cm}

\noindent \begin{tabular}{cp{10cm}}
\raisebox{-.75\height}{\includegraphics[width=0.22\textwidth]{cc_by}} &
This documentation is licensed under the Creative Commons Attribution 4.0 International
License. To view a copy of this license, visit
\href{http://creativecommons.org/licenses/by/4.0}{\color{blue}{http://creativecommons.org/licenses/by/4.0}}. \\
\end{tabular}


\vspace*{1.5cm}

{\footnotesize

\noindent Intel, the Intel logo, and other Intel marks are trademarks of Intel
Corporation in the U.S. and/or other countries. \\
\iftoggleverb{ispbs}
\noindent Altair, the Altair logo, OpenPBS, and other Altair marks are
trademarks of Altair Engineering, Inc. in the U.S. and/or other countries. \\
\fi
\noindent *Other names and brands may be claimed as the property of others. \\



}
 

\newpage
\tableofcontents
\newpage

% Introduction  ----------------------------------------------------


\section{Introduction} \label{sec:introduction}
\input{common/install_header}
\input{common/intro} \\

\input{common/base_edition/edition}
\input{common/audience}
\input{common/requirements}
% -*- mode: latex; fill-column: 120; -*- 

\subsection{Inputs} \label{sec:inputs}
As this recipe details installing a cluster starting from bare-metal, there is a requirement to define IP addresses and
gather hardware MAC addresses in order to support a controlled provisioning process. These values are necessarily unique
to the hardware being used, and this document uses variable substitution (\texttt{\$\{variable\}}) in the command-line
examples that follow to highlight where local site inputs are required. A summary of the required and optional variables
used throughout this recipe are presented below. Note that while the example definitions above correspond to a small
4-node compute subsystem, the compute parameters are defined in array format to accommodate logical extension to larger
node counts. \\

\vspace*{0.2cm}
\begin{tabular}{@{}>{\textbullet}l p{7cm} l}
& \texttt{\$\{sms\_name\}} & {\small \# Hostname for SMS server} \\
& \texttt{\$\{sms\_ip\}} & {\small \# Internal IP address on SMS server}  \\
\iftoggleverb{isxCAT}
& \texttt{\$\{domain\_name\}} & {\small \# Local network domain name}  \\
\fi
& \texttt{\$\{sms\_eth\_internal\}} & {\small \# Internal Ethernet interface on SMS} \\
\iftoggleverb{isWarewulf}
& \texttt{\$\{eth\_provision\}} & {\small \# Provisioning interface for computes} \\
\fi
& \texttt{\$\{internal\_network\}} & {\small \# IP space for internal network (e.g. 10.0.0.0)} \\
& \texttt{\$\{internal\_netmask\}} & {\small \# Subnet netmask for internal network} \\
& \texttt{\$\{ntp\_server\}} & {\small \# Local ntp server for time synchronization} \\
& \texttt{\$\{bmc\_username\}} & {\small \# BMC username for use by IPMI} \\
& \texttt{\$\{bmc\_password\}} & {\small \# BMC password for use by IPMI} \\
& \texttt{\$\{num\_computes\}} & {\small \# Total \# of desired compute nodes} \\
& \texttt{\$\{c\_ip[0]\}}, \, \texttt{\$\{c\_ip[1]\}}, ... & {\small \# Desired compute node addresses} \\
& \texttt{\$\{c\_bmc[0]\}}, \texttt{\$\{c\_bmc[1]\}}, ... & {\small \# BMC addresses for computes} \\
& \texttt{\$\{c\_mac[0]\}}, \texttt{\$\{c\_mac[1]\}}, ... & {\small \# MAC addresses for computes} \\
& \texttt{\$\{c\_name[0]\}}, \texttt{\$\{c\_name[1]\}}, ... & {\small \# Host names for computes} \\
& \texttt{\$\{compute\_regex\}} & {\small \# Regex matching all compute node names (e.g. ``c*'')} \\
& \texttt{\$\{compute\_prefix\}} & {\small \# Prefix for compute node names (e.g. ``c'')} \\
\iftoggleverb{isxCAT}
& \texttt{\$\{iso\_path\}} & {\small \# Directory location of OS iso for \xCAT{} install} \\
\nottoggle{isxCATstateful}
{& \texttt{\$\{synclist\}} & {\small \# Filesystem location of \xCAT{} synclist file} \\}
\fi
\iftoggleverb{isxCATstateful}
& \texttt{\$\{ohpc\_repo\_dir\}} & {\small \# Directory location of local \OHPC{} repository mirror} \\
& \texttt{\$\{epel\_repo\_dir\}} & {\small \# Directory location of local EPEL repository
mirror} \\
\fi
\end{tabular}

\vspace*{0.2cm}
\noindent {Optional:} 
\vspace*{0.1cm}

\begin{tabular}{@{}>{\textbullet}l p{7cm} l}
\iftoggleverb{isx86}
& \texttt{\$\{sysmgmtd\_host\}} & {\small \# BeeGFS System Management host name} \\
& \texttt{\$\{mgs\_fs\_name\}} & {\small \# Lustre MGS mount name} \\
& \texttt{\$\{sms\_ipoib\}} & {\small \# IPoIB address for SMS server} \\
& \texttt{\$\{ipoib\_netmask\}} & {\small \# Subnet netmask for internal IPoIB} \\
& \texttt{\$\{c\_ipoib[0]\}}, \texttt{\$\{c\_ipoib[1]\}}, ... & {\small \# IPoIB addresses for computes} \\
\fi
\iftoggleverb{isWarewulf}
& \texttt{\$\{kargs\}} & {\small \# Kernel boot arguments} \\
\fi
& \texttt{\$\{nagios\_web\_password\}} & {\small \# Nagios web access password} \\
\end{tabular}




% begin_ohpc_run
% ohpc_validation_newline
% ohpc_validation_comment Verify OpenHPC repository has been enabled before proceeding
% ohpc_validation_newline
% ohpc_command yum repolist | grep -q OpenHPC
% ohpc_command if [ $? -ne 0 ];then
% ohpc_command    echo "Error: OpenHPC repository must be enabled locally"
% ohpc_command    exit 1
% ohpc_command fi
% end_ohpc_run

% Base Operating System --------------------------------------------

\section{Install Base Operating System (BOS)}
\input{common/bos}

% begin_ohpc_run
% ohpc_validation_newline
% ohpc_validation_comment Disable firewall 
\begin{lstlisting}[language=bash,keywords={}]
[sms](*\#*) systemctl disable firewalld
[sms](*\#*) systemctl stop firewalld
\end{lstlisting}
% end_ohpc_run

% ------------------------------------------------------------------

\section{Install \OHPC{} Components} \label{sec:basic_install}
\input{common/install_ohpc_components_intro.tex}

\subsection{Enable \OHPC{} repository for local use} \label{sec:enable_repo}
\input{common/enable_ohpc_repo}
In addition to the \OHPC{} 
\iftoggle{isxCAT}{and \xCAT{} package repositories,}{package repository,}
the {\em master} host also requires access to the standard base OS distro
repositories in order to resolve necessary dependencies. For \baseOS{}, the
requirements are to have access to the {\color{purple}{OS}}, {\color{purple}{Everything}},
{\color{purple}{EPOL main}} and {\color{purple}{EPOL update}} repositories
for which mirrors are freely available online:

\begin{itemize*}
\item OS
  (e.g. \href{http://repo.openeuler.org/openEuler-22.03-LTS/OS/}
            {\color{blue}{http://repo.openeuler.org/openEuler-22.03-LTS/OS/}} )
\item Everything
  (e.g. \href{http://repo.openeuler.org/openEuler-22.03-LTS/everything/}
            {\color{blue}{http://repo.openeuler.org/openEuler-22.03-LTS/everything/}} )
\item EPOL-main
  (e.g. \href{http://repo.openeuler.org/openEuler-22.03-LTS/EPOL/main/}
            {\color{blue}{http://repo.openeuler.org/openEuler-22.03-LTS/EPOL/main/}} )
\item EPOL-update
  (e.g. \href{http://repo.openeuler.org/openEuler-22.03-LTS/EPOL/update/main}
            {\color{blue}{http://repo.openeuler.org/openEuler-22.03-LTS/EPOL/update/main}} )
\end{itemize*}

One might replace {\color{blue}{http://repo.openeuler.org}} with any of the mirrors listed
\href{https://www.openeuler.org/en/mirror/list/}{\color{blue}{here}}

\input{common/automation}

\subsection{Add provisioning services on {\em master} node} \label{sec:add_provisioning}
With the \OHPC{} repository enabled, we can now begin adding desired components onto the
{\em master} server. This repository provides a number of aliases that group
logical components together in order to help aid in this process. For
reference, a complete list of available group aliases and RPM packages available
via \OHPC{} are provided in Appendix~\ref{appendix:manifest}. To add
support for provisioning services, the following command adds a common base 
package followed along with the Warewulf provisioning system. Then the main 
Warewulf configuration file is edited to reflect the environment. 

%\nottoggle{isCentOS}{\clearpage}

% begin_ohpc_run
% ohpc_comment_header Add baseline OpenHPC and provisioning services \ref{sec:add_provisioning}
\begin{lstlisting}[language=bash,keywords={}]
# Install base meta-packages
[sms](*\#*) (*\install*) ohpc-base warewulf-ohpc hwloc-ohpc
\end{lstlisting}
% end_ohpc_run

\input{common/enable_pxe}
\input{common/time}

\subsection{Add resource management services on {\em master} node} \label{sec:add_rm}
\input{common/install_openpbs}

%% Add if/when IB is available for testing
%% \subsection{Optionally add \InfiniBand{} support services on {\em master} node} \label{sec:add_ofed}
%% The following command adds OFED and PSM support using base distro-provided drivers
to the chosen {\em master} host.

% begin_ohpc_run
% ohpc_comment_header Optionally add InfiniBand support services on master node \ref{sec:add_ofed}
% ohpc_command if [[ ${enable_ib} -eq 1 ]];then
% ohpc_indent 5
\begin{lstlisting}[language=bash,keywords={}]
[sms](*\#*) (*\groupinstall*) "InfiniBand Support"
# Load IB services
[sms](*\#*) udevadm trigger --type=devices --action=add
[sms](*\#*) systemctl restart rdma-load-modules@infiniband.service
\end{lstlisting}
% ohpc_indent 0
% ohpc_command fi
% end_ohpc_run

\input{common/opensm}

With the \InfiniBand{} drivers included, you can also enable (optional) IPoIB functionality
which provides a mechanism to send IP packets over the IB network. If you plan
to mount a \Lustre{} file system over \InfiniBand{} (see \S\ref{sec:lustre_client}
for additional details), then having IPoIB enabled is a requirement for the
\Lustre{} client. \OHPC{} provides a template configuration file to aid in setting up
an {\em ib0} interface on the {\em master} host. To use, copy the template
provided and update the \texttt{\$\{sms\_ipoib\}} and
\texttt{\$\{ipoib\_netmask\}} entries to match local desired settings (alter ib0
naming as appropriate if system contains dual-ported or multiple HCAs).

% begin_ohpc_run
% ohpc_validation_newline
% ohpc_validation_comment Optionally enable IPoIB interface on SMS
% ohpc_command if [[ ${enable_ipoib} -eq 1 ]];then
% ohpc_indent 5
% ohpc_validation_comment Enable ib0
\begin{lstlisting}[language=bash,literate={-}{-}1,keywords={},upquote=true]
[sms](*\#*) cp /opt/ohpc/pub/examples/network/centos/ifcfg-ib0 /etc/sysconfig/network-scripts

# Define local IPoIB address and netmask
[sms](*\#*) perl -pi -e "s/master_ipoib/${sms_ipoib}/" /etc/sysconfig/network-scripts/ifcfg-ib0
[sms](*\#*) perl -pi -e "s/ipoib_netmask/${ipoib_netmask}/" /etc/sysconfig/network-scripts/ifcfg-ib0

# configure NetworkManager to *not* override local /etc/resolv.conf
[sms](*\#*) echo "[main]"   >  /etc/NetworkManager/conf.d/90-dns-none.conf
[sms](*\#*) echo "dns=none" >> /etc/NetworkManager/conf.d/90-dns-none.conf
# Start up NetworkManager to initiate ib0
[sms](*\#*) systemctl start NetworkManager

\end{lstlisting}
% ohpc_indent 0
% ohpc_command fi
% end_ohpc_run



%\vspace*{-0.15cm}
\subsection{Complete basic Warewulf setup for {\em master} node} \label{sec:setup_ww}
At this point, all of the packages necessary to use \Warewulf{} on the {\em
master} host should be installed. Next, we need to update the configuration 
to allow \Warewulf{} to work with \baseOS{} and to
support local provisioning using a second private interface (refer to
Figure~\ref{fig:physical_arch}).
%\vspace*{-0.05cm}

% begin_ohpc_run
% ohpc_comment_header Complete basic Warewulf setup for master node \ref{sec:setup_ww}
%\begin{verbatim}

\begin{lstlisting}[language=bash,literate={-}{-}1,keywords={},upquote=true,keepspaces]
# Enable internal interface for provisioning
[sms](*\#*) ip link set dev ${sms_eth_internal} up
[sms](*\#*) ip address add ${sms_ip}/${internal_netmask} broadcast + dev ${sms_eth_internal}

# Edit the warewulf.conf file to use appropriate interface and settings
[sms](*\#*) sed -i "s/ipaddr:.*/ipaddr: ${sms_ip}/g" /etc/warewulf/warewulf.conf
[sms](*\#*) sed -i "s/netmask:.*/netmask: ${internal_netmask}/g" /etc/warewulf/warewulf.conf
[sms](*\#*) sed -i "s/network:.*/network: ${internal_network}/g" /etc/warewulf/warewulf.conf
[sms](*\#*) sed -i "s/range start:.*/range start: ${c_ip[0]}/g" /etc/warewulf/warewulf.conf
[sms](*\#*) sed -i "s/range start:.*/range start: ${c_ip[$num_computes]}/g" /etc/warewulf/warewulf.conf
[sms](*\#*) sed -i "s/mount: false/mount: true/g" /etc/warewulf/warewulf.conf

# Configuring Warewulf will restart/enable relevant services to support provisioning
[sms](*\#*) systemctl enable --now warewulfd
[sms](*\#*) wwctl configure --all

# This line is a fix for a problem with the current Warewulf rpm
# Hopefully user's can skip this step in a future release
[sms](*\#*) cp -r /srv/tftpboot/warewulf /var/lib/tftpboot/ 
\end{lstlisting}
%\end{verbatim}
% end_ohpc_run


\subsection{Define {\em compute} image for provisioning}
With the provisioning services enabled, the next step is to define and
customize a system image that can subsequently be used to provision one or more
{\em compute} nodes. The following subsections highlight this process.

\subsubsection{Build initial BOS image} \label{sec:assemble_bos}
The \OHPC{} build of \Warewulf{} includes specific enhancements enabling support for
\baseOS{}. The following steps illustrate the process to build a minimal, default
image for use with \Warewulf{}. We begin by defining a directory structure on the 
{\em master} host that will represent the root filesystem of the compute node. The 
default location for this example is in
\texttt{/opt/ohpc/admin/images/\baseos{}}.

\begin{center}
  \begin{tcolorbox}[]
    \small \Warewulf{} is configured by default to access an external
    repository (repo.openeuler.org) during the \texttt{wwmkchroot} process.  If
    the master host cannot reach the public openEuler mirrors, or if you prefer to
    access a locally cached mirror, set the \texttt{\$\{YUM\_MIRROR\}}
    environment variable to your desired repo location {\em prior} to running
    the \texttt{wwmkchroot} command below. For example:

% begin_ohpc_run
% ohpc_command if [ ! -z ${BOS_MIRROR+x} ]; then
% ohpc_indent 5
\begin{lstlisting}[language=bash,keywords={}]
# Override default OS repository (optional) - set YUM_MIRROR variable to desired repo location
[sms](*\#*) export YUM_MIRROR=${BOS_MIRROR}
\end{lstlisting}
% ohpc_indent 0
% ohpc_command fi
% end_ohpc_runƒ

\end{tcolorbox}
\end{center}

% begin_ohpc_run
% ohpc_comment_header Create compute image for Warewulf \ref{sec:assemble_bos}
\begin{lstlisting}[language=bash,literate={-}{-}1,keywords={},upquote=true,keepspaces,literate={BOSVER}{\baseos{}}1]
# Define chroot location 
[sms](*\#*) export CHROOT=/opt/ohpc/admin/images/BOSVER

# Build initial chroot image
[sms](*\#*) wwmkchroot -v openeuler-22.03 $CHROOT
# Enable OpenHPC and openEuler repos inside chroot
[sms](*\#*) dnf -y --installroot $CHROOT install openEuler-release
[sms](*\#*) cp -p /etc/yum.repos.d/OpenHPC*.repo $CHROOT/etc/yum.repos.d
\end{lstlisting}
% end_ohpc_run


\subsubsection{Add \OHPC{} components} \label{sec:add_components}
The process used in the previous step is designed to
provide a minimal \baseOS{} configuration. Next, we add additional components
to include resource management client services, NTP support, and
other additional packages to support the default \OHPC{} environment. This
process augments the chroot-based install performed by \texttt{wwmkchroot} to
modify the base provisioning image and will access the BOS and \OHPC{}
repositories to resolve package install requests. We begin by installing a few
common base packages:

% begin_ohpc_run
% ohpc_comment_header Add OpenHPC base components to compute image \ref{sec:add_components}
\begin{lstlisting}[language=bash,literate={-}{-}1,keywords={},upquote=true]
# Install compute node base meta-package
[sms](*\#*) (*\chrootinstall*) ohpc-base-compute
\end{lstlisting}
% end_ohpc_run

\noindent Now, we can include additional required components to the compute
instance including resource manager client, NTP, and development environment modules support.



% begin_ohpc_run
% ohpc_validation_comment Add OpenHPC components to compute instance
\begin{lstlisting}[language=bash,literate={-}{-}1,keywords={},upquote=true]
# Add OpenPBS client support
[sms](*\#*) (*\chrootinstall*) openpbs-execution-ohpc
[sms](*\#*) perl -pi -e "s/PBS_SERVER=\S+/PBS_SERVER=${sms_name}/" $CHROOT/etc/pbs.conf
[sms](*\#*) echo "PBS_LEAF_NAME=${sms_name}" >> /etc/pbs.conf
[sms](*\#*) chroot $CHROOT opt/pbs/libexec/pbs_habitat
[sms](*\#*) perl -pi -e "s/\$clienthost \S+/\$clienthost ${sms_name}/" $CHROOT/var/spool/pbs/mom_priv/config
[sms](*\#*) echo "\$usecp *:/home /home" >> $CHROOT/var/spool/pbs/mom_priv/config
[sms](*\#*) chroot $CHROOT systemctl enable pbs

# Add Network Time Protocol (NTP) support
[sms](*\#*) (*\chrootinstall*) chrony
[sms](*\#*) echo "server ${sms_ip} iburst" >> $CHROOT/etc/chrony.conf

# Add kernel drivers (matching kernel version on SMS node)
[sms](*\#*) (*\chrootinstall*) kernel-`uname -r`

# Include modules user environment
[sms](*\#*) (*\chrootinstall*) lmod-ohpc
\end{lstlisting}
% end_ohpc_run

\subsubsection{Customize system configuration} \label{sec:master_customization}
Prior to assembling the image, it is advantageous to perform any additional
customization within the chroot environment created for the desired {\em
compute} instance. Older versions of \Warewulf{} required manual steps to
create and share ssh keys supporting local access and configure NFS mounts
from the {\em master} node, but \Warewulf{} 4 handles these things automatically 
through the configuration file. 

\input{common/restart_nfs}

%\clearpage
\subsubsection{Additional Customization ({\em optional})} \label{sec:addl_customizations}
\input{common/compute_customizations_intro}

%% Add if/when IB is available for testing
%% \paragraph{Increase locked memory limits}
%% \input{common/memlimits}

%%\paragraph{Add \Lustre{} client} \label{sec:lustre_client}
%%\input{common/lustre-client}
%%\input{common/lustre-client-centos}
%%\vspace*{0.5cm}
%%\input{common/lustre-client-post}

%\clearpage
\vspace*{-.2cm}
\paragraph{Enable forwarding of system logs} \label{sec:add_syslog}
\input{common/syslog}

%\vspace*{-.2cm}
\clearpage
\paragraph{Add \Nagios{} monitoring} \label{sec:add_nagios}
\input{common/nagios}

\clearpage
%\vspace*{-0.32cm}
\paragraph{Add \clustershell{}}
\input{common/clustershell}

%\clearpage
%\paragraph{Add \mrsh{}}
%\input{common/mrsh}

\paragraph{Add \genders{}}
\input{common/genders}

\vspace*{-.1cm}
\paragraph{Add Magpie}
\input{common/magpie}

\vspace*{-.1cm}
\paragraph{Add \conman{}} \label{sec:add_conman}
\conman{} is a serial console management program designed to support a large
number of console devices and simultaneous users. It supports logging console
device output and connecting to compute node consoles via IPMI
serial-over-lan. Installation and example configuration is outlined below.

% begin_ohpc_run
% ohpc_validation_newline
% ohpc_validation_comment Optionally, enable conman and configure
% ohpc_command if [[ ${enable_ipmisol} -eq 1 ]];then
% ohpc_indent 5
\begin{lstlisting}[language=bash,keywords={},upquote=true]
# Install conman to provide a front-end to compute consoles and log output
[sms](*\#*) (*\install*) conman-ohpc

# Configure conman for computes (note your IPMI password is required for console access)
[sms](*\#*) for ((i=0; i<$num_computes; i++)) ; do
              echo -n 'CONSOLE name="'${c_name[$i]}'" dev="ipmi:'${c_bmc[$i]}'" '
              echo 'ipmiopts="'U:${bmc_username},P:${IPMI_PASSWORD:-undefined},W:solpayloadsize'"'
        done >> /etc/conman.conf

# Enable and start conman
[sms](*\#*) systemctl enable conman
[sms](*\#*) systemctl start conman
\end{lstlisting}
% ohpc_indent 0
% ohpc_command fi
% end_ohpc_run

\iftoggleverb{isWarewulf}
\noindent Note that an additional kernel boot option is typically necessary to
enable serial console output. This option is highlighted after
compute nodes have been registered with the provisioning system.
\fi

\iftoggleverb{isxCAT}
\noindent Note that additional options are typically necessary to
enable serial console output. These are setup during the node registration
process in \S\ref{sec:xcat_add_nodes} 
\fi




\vspace*{-.1cm}
\paragraph{Add \nhc{}} \label{sec:add_nhc}
\input{common/nhc}

\subsubsection{Import files} \label{sec:file_import}
The \Warewulf{} system includes functionality to import arbitrary files from
the provisioning server for distribution to managed hosts through a system 
called "overlays". Some files, like \texttt{/etc/passwd}, and \texttt{/etc/hosts}
handled in this way by default.  Here we add directories and files to the 
\texttt{generic} overlay that is applied to all nodes. 

% begin_ohpc_run
% ohpc_comment_header Import files \ref{sec:file_import}
\begin{lstlisting}[language=bash,literate={-}{-}1,keywords={},upquote=true]
[sms](*\#*) wwctl overlay edit --parents generic /etc/sysconfig/slurmd
# in the editor, add the following line (without the preceding #)
# SLURMD_OPTIONS="--conf-server 10.0.0.3"

[sms](*\#*) wwctl overlay mkdir generic --mode 0700 /etc/munge
[sms](*\#*) wwctl overlay import generic /etc/munge/munge.key
[sms](*\#*) wwctl overlay chown generic /etc/munge/munge.key $(id -u munge) $(id -g munge)
[sms](*\#*) wwctl overlay chown generic /etc/munge $(id -u munge) $(id -g munge)
[sms](*\#*) wwctl overlay edit --parents generic /etc/security/access.conf
# in the editor, add the following lines (without the preceding #'s)
# + : root : ALL
# + : ssh-allowed : ALL
# - : ALL : ALL

[sms](*\#*) wwctl overlay mkdir generic /etc/pam.d
[sms](*\#*) wwctl overlay import generic $CHROOT/etc/pam.d/sshd /etc/pam.d/sshd
[sms](*\#*) wwctl overlay edit generic /etc/pam.d/sshd
# in the editor, add the following lines (without the preceding #'s)
# account sufficient pam_access.so
# account required pam_slurm.so

# add the following to support unprivileged user namespaces for tools like Apptainer
[sms](*\#*) wwctl overlay import generic /etc/subuid
[sms](*\#*) wwctl overlay import generic /etc/subgid

# there is a bug in the current ohpc-supported version of Warewulf that produces
# improper ownership. these lines will fix.
[sms](*\#*) chown munge:munge /srv/warewulf/overlays/generic/etc/munge/munge.key
[sms](*\#*) chown munge:munge /srv/warewulf/overlays/generic/etc/munge/

\end{lstlisting}
% \end_ohpc_run

%% \input{common/import_ww_files_ib_centos}
\subsection{Finalizing provisioning configuration} \label{sec:assemble_bootstrap}

\Warewulf{} employs a two-stage boot process for provisioning nodes via
creation of a bootstrap image that is used to initialize the process, and a virtual node
file system capsule containing the full system image. This section highlights
creation of the necessary provisioning images, followed by the registration of
desired compute nodes.

\subsubsection{Assemble bootstrap image}

The bootstrap image includes the runtime kernel and associated modules, as well
as some simple scripts to complete the provisioning process.

%\iftoggle{isCentOS_ww_slurm_aarch}{\clearpage}

% begin_ohpc_run
% ohpc_comment_header Assemble bootstrap image \ref{sec:assemble_bootstrap}
\begin{lstlisting}[language=bash,literate={-}{-}1,keywords={},upquote=true]
# Build bootstrap image
[sms](*\#*) wwctl container build rocky-8.9
\end{lstlisting}
% end_ohpc_run

\iftoggle{isCentOS_ww_slurm_aarch}{\vspace*{0.4cm}}

\iftoggle{isSLES_ww_slurm_aarch}{\vspace*{-0.1cm}}

\subsubsection{Register nodes for provisioning}

Nodes can be registered for provisioning using the following syntax.

\vspace*{0.2cm}
%\iftoggle{isx86}{\clearpage}
% begin_ohpc_run
% ohpc_validation_comment Add hosts to cluster

\begin{lstlisting}[language=bash,keywords={},upquote=true,basicstyle=\footnotesize\ttfamily,]
# first add the nodes with --discoverable so warewulf will identify new mac addresses
[sms](*\#*) wwctl node add ${compute_prefix}[1-${num_computes}] --discoverable --ipaddr ${c_ip[0]}

# next set the default profile for all nodes to contain the image we created
[sms](*\#*) wwctl profile set --yes --container rocky-8.9 "default"

# finally, build the overlays for all the nodes
[sms](*\#*) wwct overlay build
\end{lstlisting}
% end_ohpc_run

\end{lstlisting}


\input{common/add_ww_hosts_pbs}
\input{common/add_ww_hosts_finalize}

\vspace*{-0.2cm}
\subsubsection{Optional kernel arguments} \label{sec:optional_kargs}
\input{common/conman_post}
\input{common/kargs_post}

\vspace*{-0.2cm}
\subsubsection{Optionally configure stateful provisioning}
\input{common/stateful}

\vspace*{-0.1cm}
\subsection{Boot compute nodes} \label{sec:boot_computes}
At this point, the {\em master} server should be able to boot the newly defined
compute nodes. Assuming that the compute node BIOS settings are configured to
boot over PXE, all that is required to initiate the provisioning process is to
power cycle each of the desired hosts using IPMI access.
The following commands use the \texttt{ipmitool} utility to initiate power
resets on each of the four compute hosts. Note that the utility requires that
the \texttt{IPMI\_PASSWORD} environment variable be set with the local BMC password in
order to work interactively.

% begin_ohpc_run
% ohpc_comment_header Boot compute nodes \ref{sec:boot_computes}
\begin{lstlisting}[language=bash,keywords={},upquote=true]
[sms](*\#*) for ((i=0; i<${num_computes}; i++)) ; do
             ipmitool -E -I lanplus -H ${c_bmc[$i]} -U ${bmc_username} -P ${bmc_password} chassis power reset
        done
\end{lstlisting} 
% end_ohpc_run

Once kicked off, the boot process should take less than 5 minutes (depending on
BIOS post times) and you can verify that the compute hosts are available via
ssh, or via parallel ssh tools to multiple hosts. For example, to run a command
on the newly imaged compute hosts using \texttt{pdsh}, execute the following:

%\iftoggle{isCentOS_ww_pbs_x86}{\clearpage}
%\iftoggle{isCentOS}{\clearpage}
  
\begin{lstlisting}[language=bash]
[sms](*\#*) pdsh -w c[1-4] uptime
c1  05:03am  up   0:02,  0 users,  load average: 0.20, 0.13, 0.05
c2  05:03am  up   0:02,  0 users,  load average: 0.20, 0.14, 0.06
c3  05:03am  up   0:02,  0 users,  load average: 0.19, 0.15, 0.06
c4  05:03am  up   0:02,  0 users,  load average: 0.15, 0.12, 0.05
\end{lstlisting}

 

\vspace*{-0.50cm}
\section{Install \OHPC{} Development Components} \label{sec:install_dev}
\input{common/dev_intro.tex}

\vspace*{-0.15cm}
\subsection{Development Tools} \label{sec:install_dev_tools}
\input{common/dev_tools}

\vspace*{-0.15cm}
\subsection{Compilers} \label{sec:install_compilers}
\OHPC{} presently packages the \GNU{} compiler toolchain integrated with the 
underlying Lmod modules system in a hierarchical fashion. The modules
system will conditionally present compiler-dependent software based on the
toolchain currently loaded. 

% begin_ohpc_run
% ohpc_comment_header Install Compilers \ref{sec:install_compilers}
\begin{lstlisting}[language=bash]
[sms](*\#*) (*\install*) gnu12-compilers-ohpc
\end{lstlisting}
% end_ohpc_run

%%% The llvm compiler toolchains are also provided as a standalone additional
%%% compiler family (ie. users can easily switch between gcc/clang environments),
%%% but we do not provide the full complement of downstream library builds.
%%% 
%%% % begin_ohpc_run
%%% % ohpc_comment_header Install llvm Compilers
%%% \begin{lstlisting}[language=bash]
%%% [sms](*\#*) (*\install*) llvm5-compilers-ohpc
%%% \end{lstlisting}
%%% % end_ohpc_run


%\clearpage
\vspace*{0.5cm}
\subsection{MPI Stacks} \label{sec:mpi}
For MPI development and runtime support, \OHPC{} provides pre-packaged builds
for two MPI families that are compatible with both ethernet and high-speed
fabrics.  These MPI stacks can be installed as follows:

% begin_ohpc_run
% ohpc_comment_header Install MPI Stacks \ref{sec:mpi}
% ohpc_command if [[ ${enable_mpi_defaults} -eq 1 ]];then
% ohpc_indent 5
\begin{lstlisting}[language=bash]
[sms](*\#*) (*\install*) openmpi4-gnu12-ohpc mpich-ofi-gnu12-ohpc
\end{lstlisting}
% ohpc_indent 0
% ohpc_command fi
% end_ohpc_run

Note that OpenHPC 2.x introduces the use of two related transport layers for
the MPICH and OpenMPI builds that support a variety of underlying
fabrics: \href{https://www.openucx.org}{UCX} (Unified Communication X)
and \href{https://ofiwg.github.io/libfabric/}{OFI} (OpenFabrics interfaces).
In the case of OpenMPI, a monolithic build is provided which supports both
transports and end-users can customize their runtime preferences with
environment variables. For MPICH, two separate builds are provided and the
example above highlighted installing the {\texttt ofi} variant.  However, the
packaging is designed such that both versions can be installed simultaneously
and users can switch between the two via normal module command
semantics. Alternatively, a site can choose to install the {\texttt ucx} variant
instead as a drop-in MPICH replacement:

% begin_ohpc_run
% ohpc_command if [[ ${enable_mpich_ucx} -eq 1 ]];then
% ohpc_indent 5
\begin{lstlisting}[language=bash]
[sms](*\#*) (*\install*) mpich-ucx-gnu12-ohpc
\end{lstlisting}
% ohpc_indent 0
% ohpc_command fi
% end_ohpc_run

In the case where both MPICH variants are installed, two modules will be
visible in the end-user environment and an example of this configuration is
highlighted is below. 

\begin{lstlisting}[language=bash]
[sms](*\#*) module avail mpich

-------------------- /opt/ohpc/pub/moduledeps/gnu12 ---------------------
   mpich/3.4.3-ofi    mpich/3.4.3-ucx (D)
\end{lstlisting}







\subsection{Performance Tools} \label{sec:install_perf_tools}
\OHPC{} provides a variety of open-source tools to aid in application 
performance analysis (refer to Appendix~\ref{appendix:manifest} for a listing
of available packages). This group of tools can be installed as follows:

% begin_ohpc_run
% ohpc_comment_header Install Performance Tools \ref{sec:install_perf_tools}
\begin{lstlisting}[language=bash,keywords={},literate={-}{-}1]
# Install perf-tools meta-package
[sms](*\#*) (*\install*) ohpc-gnu12-perf-tools
\end{lstlisting}
% end_ohpc_run


\subsection{Setup default development environment}
System users often find it convenient to have a default development environment
in place so that compilation can be performed directly for parallel programs
requiring MPI. This setup can be conveniently enabled via modules and the \OHPC{}
modules environment is pre-configured to load an \texttt{ohpc} module on login
(if present). The following package install provides a default
environment that enables autotools, the \GNU{} compiler toolchain, and the
OpenMPI stack.

% begin_ohpc_run
\begin{lstlisting}[language=bash]
[sms](*\#*) (*\install*) lmod-defaults-gnu12-openmpi4-ohpc
\end{lstlisting}
% end_ohpc_run

\begin{center}
\begin{tcolorbox}[]
\small
\iftoggleverb{isx86}
If you want to change the default environment from the suggestion above, \OHPC{}
also provides the \GNU{} compiler toolchain with the MPICH and MVAPICH2 stacks:
\fi

\iftoggleverb{isaarch}
If you want to change the default environment from the suggestion above, \OHPC{}
also provides additional default options using the \GNU{} compiler toolchain
with multiple MPICH variants or MVAPICH2. Relevant lmod-default packages names
are as follows:
\fi

\begin{itemize*}
\item lmod-defaults-gnu12-mpich-ofi-ohpc
\item lmod-defaults-gnu12-mpich-ucx-ohpc
\iftoggleverb{isx86}
\item lmod-defaults-gnu12-mvapich2-ohpc
\fi
\end{itemize*}
\end{tcolorbox}
\end{center}


%\clearpage
\subsection{3rd Party Libraries and Tools} \label{sec:3rdparty}
\input{common/third_party_libs_intro}
\iftoggleverb{isx86}
\begin{lstlisting}[language=bash,keywords={}]
[sms](*\#*) yum search petsc-gnu12 ohpc
Loaded plugins: fastestmirror
Loading mirror speeds from cached hostfile
=========================== N/S matched: petsc-gnu9, ohpc ===========================
petsc-gnu12-impi-ohpc.x86_64 : Portable Extensible Toolkit for Scientific Computation
petsc-gnu12-mpich-ohpc.x86_64 : Portable Extensible Toolkit for Scientific Computation
petsc-gnu12-mvapich2-ohpc.x86_64 : Portable Extensible Toolkit for Scientific Computation
petsc-gnu12-openmpi4-ohpc.x86_64 : Portable Extensible Toolkit for Scientific Computation
\end{lstlisting}
\fi

\iftoggleverb{isaarch}
\begin{lstlisting}[language=bash,keywords={}]
[sms](*\#*) yum search petsc-gnu12 ohpc
Loaded plugins: fastestmirror
Loading mirror speeds from cached hostfile
=========================== N/S matched: petsc-gnu9, ohpc ===========================
petsc-gnu12-mpich-ohpc.x86_64 : Portable Extensible Toolkit for Scientific Computation
petsc-gnu12-openmpi4-ohpc.x86_64 : Portable Extensible Toolkit for Scientific Computation
\end{lstlisting}
\fi


\vspace*{.4cm}
\begin{center}
\begin{tcolorbox}[]
\small
\OHPC{}-provided 3rd party builds are configured to be installed
into a common top-level repository so that they can be easily exported to
desired hosts within the cluster. This common top-level path
(\path{/opt/ohpc/pub}) was previously configured to be mounted on {\em
 compute} nodes in \S\ref{sec:master_customization}, so the packages will be
immediately available for use on the cluster after installation on the {\em
 master} host.
\end{tcolorbox}
\end{center}

%\iftoggle{isCentOS}{\clearpage}
%\nottoggle{isCentOS}{\clearpage}

For convenience, \OHPC{} provides package aliases for these 3rd party libraries
and utilities that can be used to install available libraries for use with the
GNU compiler family toolchain. For parallel libraries, aliases are grouped by
MPI family toolchain so that administrators can choose a subset should they
favor a particular MPI stack.  Please refer to Appendix~\ref{appendix:manifest}
for a more detailed listing of all available packages in each of these functional
areas. To install all available package offerings within \OHPC{}, issue the
following:

% begin_ohpc_run
% ohpc_comment_header Install 3rd Party Libraries and Tools \ref{sec:3rdparty}
\begin{lstlisting}[language=bash,keywords={},upquote=true,keepspaces]
# Install 3rd party libraries/tools meta-packages built with GNU toolchain
[sms](*\#*) (*\install*) ohpc-gnu12-serial-libs
[sms](*\#*) (*\install*) ohpc-gnu12-io-libs
[sms](*\#*) (*\install*) ohpc-gnu12-python-libs
[sms](*\#*) (*\install*) ohpc-gnu12-runtimes
\end{lstlisting}
% end_ohpc_run





% begin_ohpc_run
% ohpc_comment_header Install 3rd Party Libraries and Tools \ref{sec:3rdparty}
% ohpc_command if [[ ${enable_mpi_defaults} -eq 1 ]];then
% ohpc_indent 5
\begin{lstlisting}[language=bash,keywords={},upquote=true,keepspaces]
# Install parallel lib meta-packages for all available MPI toolchains
[sms](*\#*) (*\install*) ohpc-gnu12-mpich-parallel-libs
[sms](*\#*) (*\install*) ohpc-gnu12-openmpi4-parallel-libs
\end{lstlisting}
% ohpc_indent 0
% ohpc_command fi
% end_ohpc_run


\subsection{Optional Development Tool Builds} \label{sec:3rdparty_arm}
In addition to the 3rd party development libraries built using the open source
toolchains mentioned in \S\ref{sec:3rdparty}, \OHPC{} also provides a subset of
{\em optional} builds compatible with the Arm Compiler for Linux.
These packages provide a similar hierarchical user environment experience as
other compiler families present in \OHPC{}.  To take advantage of the
available builds, the version 22.1 of the Arm Compiler for Linux (and any
required licenses) must be downloaded and installed separately. See the
following for more information on obtaining this toolchain:

\begin{center}
  \href{https://developer.arm.com/downloads/-/arm-compiler-for-linux}
       {\color{blue}{https://developer.arm.com/downloads/-/arm-compiler-for-linux}}
\end{center}



\begin{center}
\begin{tcolorbox}[]
As noted in \S\ref{sec:master_customization}, the default installation path for
\OHPC{} (\texttt{/opt/ohpc/pub}) is exported over NFS from the {\em master} to the 
compute nodes, but the Arm Linux compiler installer defaults to a path of 
\texttt{/opt/arm}. To make the Arm Linux compiler available to the compute 
nodes one must add an additional NFS export
for \texttt{/opt/arm} that is mounted on desired compute nodes.
\end{tcolorbox}
\end{center}

Once installed locally, the \OHPC{}
compatible packages can be installed using standard package manager semantics
after installation of a compatibility package.
\noindent To enable all 3rd party builds available in \OHPC{} that are compatible with
Arm Linux Compiler, issue the following:

% begin_ohpc_run
% ohpc_comment_header Install Optional Development Tools for use with Arm Linux Compiler \ref{sec:3rdparty_arm}
% ohpc_command if [[ ${enable_arm1_packages} -eq 1 ]];then
% ohpc_indent 5
\begin{lstlisting}[language=bash,keywords={},upquote=true,keepspaces]
# Install OpenHPC compatibility packages (requires prior installation of Arm Linux Compiler)
[sms](*\#*) (*\install*) arm1-compilers-devel-ohpc
\end{lstlisting}

\iftoggle{isSLES_ww_slurm_x86}{\clearpage}

\begin{lstlisting}[language=bash,keywords={},upquote=true,keepspaces]
# Install 3rd party libraries/tools meta-packages built with Arm vendor toolchain
[sms](*\#*) (*\install*) ohpc-arm1-serial-libs
[sms](*\#*) (*\install*) ohpc-arm1-io-libs
[sms](*\#*) (*\install*) ohpc-arm1-perf-tools
[sms](*\#*) (*\install*) ohpc-arm1-python3-tools

[sms](*\#*) (*\install*) ohpc-arm1-mpich-parallel-libs
[sms](*\#*) (*\install*) ohpc-arm1-openmpi4-parallel-libs
\end{lstlisting}
% ohpc_indent 0
% ohpc_command fi
% end_ohpc_run

% add back when ready
%[sms](*\#*) (*\install*) ohpc-arm1-python-libs



\clearpage
\section{Resource Manager Startup} \label{sec:rms_startup}
\input{common/openpbs_startup}

\section{Run a Test Job} \label{sec:test_job}
With the resource manager enabled for production usage, users should now be
able to run jobs. To demonstrate this, we will add a ``test'' user on the {\em master}
host that can be used to run an example job.

% begin_ohpc_run
\begin{lstlisting}[language=bash,keywords={}]
[sms](*\#*) useradd -m test
\end{lstlisting}
% end_ohpc_run

\Warewulf{} installs a utility on the compute nodes to automatically 
synchronize known files from the provisioning server at five minute intervals. In this
recipe, recall that we previously registered credential files with Warewulf (e.g. passwd,
group, and shadow) so that these files would be propagated during compute node
imaging. However, with the addition of a new ``test'' user above, the files
have been outdated and we need to update the Warewulf database to incorporate
the additions. This re-sync process can be accomplished as follows:

% begin_ohpc_run
\begin{lstlisting}[language=bash,keywords={}]
[sms](*\#*) wwsh file resync passwd shadow group
\end{lstlisting}
% end_ohpc_run

% begin_ohpc_run
% ohpc_command sleep 2
% end_ohpc_run

\begin{center}
\begin{tcolorbox}[]
\small
After re-syncing to notify Warewulf of file modifications made on the {\em
master} host, it should take approximately 5 minutes for the changes to
propagate. However, you can also manually pull the changes from compute nodes
via the following:
% begin_ohpc_run
\begin{lstlisting}[language=bash,keywords={}]
[sms](*\#*) pdsh -w $compute_prefix[1-${num_computes}] /warewulf/bin/wwgetfiles
\end{lstlisting}
% end_ohpc_run
\end{tcolorbox}
\end{center}

\input{common/prun}

%\iftoggle{isSLES_ww_pbs_x86}{\clearpage}

\subsection{Interactive execution}
To use the newly created ``test'' account to compile and execute the
application {\em interactively} through the resource manager, execute the
following (note the use of \texttt{mpiexec} for parallel job launch which summarizes
the underlying native job launch mechanism being used):

\begin{lstlisting}[language=bash,keywords={}]
# Switch to "test" user
[sms](*\#*) su - test

# Compile MPI "hello world" example
[test@sms ~]$ mpicc -O3 /opt/ohpc/pub/examples/mpi/hello.c

# Submit interactive job request and use prun to launch executable
[test@sms ~]$ qsub -I -l select=2:mpiprocs=4

[test@c1 ~]$ prun ./a.out

[prun] Master compute host = c1
[prun] Resource manager = openpbs
[prun] Launch cmd = mpiexec.hydra -rmk pbs ./a.out

 Hello, world (8 procs total)
     --> Process #   0 of   8 is alive. -> c1
     --> Process #   1 of   8 is alive. -> c1
     --> Process #   2 of   8 is alive. -> c1
     --> Process #   3 of   8 is alive. -> c1
     --> Process #   4 of   8 is alive. -> c2
     --> Process #   5 of   8 is alive. -> c2
     --> Process #   6 of   8 is alive. -> c2
     --> Process #   7 of   8 is alive. -> c2
\end{lstlisting}


\begin{center}
\begin{tcolorbox}[]
The following table provides approximate command equivalences between OpenPBS
and SLURM:

\vspace*{0.15cm}
\input common/rms_equivalence_table
\end{tcolorbox}
\end{center}

%\iftoggle{isCentOS_ww_pbs_x86}{\clearpage}

\iftoggle{isSLES_ww_pbs_x86}{\clearpage}

\iftoggle{isSLES_ww_pbs_aarch}{\clearpage}

\subsection{Batch execution}

For batch execution, \OHPC{} provides a simple job script for reference (also
housed in the \path{/opt/ohpc/pub/examples} directory. This example script can
be used as a starting point for submitting batch jobs to the resource manager
and the example below illustrates use of the script to submit a batch job for
execution using the same executable referenced in the previous interactive
example.

\begin{lstlisting}[language=bash,keywords={}]
# Copy example job script
[test@sms ~]$ cp /opt/ohpc/pub/examples/openpbs/job.mpi .

# Examine contents (and edit to set desired job sizing characteristics)
[test@sms ~]$ cat job.mpi
#!/bin/bash
#----------------------------------------------------------
# Job name
#PBS -N test

# Name of stdout output file
#PBS -o job.out

# Total number of nodes and MPI tasks/node requested
#PBS -l select=2:mpiprocs=4

# Run time (hh:mm:ss) - 1.5 hours
#PBS -l walltime=01:30:00
#----------------------------------------------------------

# Change to submission directory
cd $PBS_O_WORKDIR

# Launch MPI-based executable
prun ./a.out

# Submit job for batch execution
[test@sms ~]$ qsub job.mpi
5.master

\end{lstlisting}


\clearpage
\appendix
{\bf \LARGE \centerline{Appendices}} \vspace*{0.2cm}

\addcontentsline{toc}{section}{Appendices}
\renewcommand{\thesubsection}{\Alph{subsection}}

\input{common/automation_appendix}
\subsection{Upgrading OpenHPC Packages}  \label{appendix:upgrade}


As newer \OHPC{} releases are made available, users are encouraged to upgrade
their locally installed packages against the latest repository versions to
obtain access to bug fixes and newer component versions. This can be
accomplished with the underlying package manager as \OHPC{} packaging maintains
versioning state across releases. Also, package builds available from the
\OHPC{} repositories have ``\texttt{-ohpc}'' appended to their names so that
wild cards can be used as a simple way to obtain updates. The following general
procedure highlights a method for upgrading existing installations.
When upgrading from a minor release older than v\OHPCVerTree{}, you will first
need to update your local \OHPC{} repository configuration to point against the
v\OHPCVerTree{} release (or update your locally hosted mirror). Refer to
\S\ref{sec:enable_repo} for more details on enabling the latest
repository. In contrast, when upgrading between micro releases on the same
branch (e.g. from v\OHPCVerTree{} to \OHPCVerTree{}.2), there is no need to
adjust local package manager configurations when using the public repository as
rolling updates are pre-configured.
 
\begin{enumerate*}
\item (Optional) Ensure repo metadata is current (on head node and in chroot
  location(s)). Package managers will naturally do this on their own over time,
  but if you are wanting to access updates immediately after a new release,
  the following can be used to sync to the latest.

\begin{lstlisting}[language=bash,keywords={}]
[sms](*\#*)  (*\clean*)
[sms](*\#*)  (*\chrootclean*)
\end{lstlisting}

\item Upgrade master (SMS) node

\begin{lstlisting}[language=bash,keywords={}]
[sms](*\#*)  (*\upgrade*) "*-ohpc"

# Any new Base OS provided dependencies can be installed by
# updating the ohpc-base metapackage
[sms](*\#*)  (*\upgrade*) "ohpc-base"
\end{lstlisting}
  
\item Upgrade packages in compute image

\begin{lstlisting}[language=bash,keywords={}]
[sms](*\#*)  (*\chrootupgrade*) "*-ohpc"

# Any new compute-node Base OS provided dependencies can be installed by
# updating the ohpc-base-compute metapackage
[sms](*\#*)  (*\chrootupgrade*) "ohpc-base-compute"
\end{lstlisting}
  
\item Rebuild image(s)

\iftoggleverb{isWarewulf}
\begin{lstlisting}[language=bash,keywords={}]
[sms](*\#*) wwvnfs --chroot $CHROOT
\end{lstlisting}
\fi

\iftoggleverb{isxCAT}
\begin{lstlisting}[language=bash,keywords={},basicstyle=\fontencoding{T1}\fontsize{8.0}{10}\ttfamily,
    literate={-}{-}1 {BOSVER}{\baseos{}}1 {ARCH}{\arch{}}1]
[sms](*\#*) packimage BOSVER-x86_64-netboot-compute
\end{lstlisting}
\fi

\end{enumerate*}

\noindent In the case where packages were upgraded within the chroot compute image,
you will need to reboot the compute nodes when convenient to enable the
changes.

%%% \subsubsection{New component variants}
%%% 
%%% As newer variants of key compiler/MPI stacks are released, \OHPC{} will
%%% periodically add toolchains enabling the latest variant. To stay consistent
%%% throughout the build hierarchy, minimize recompilation requirements for existing
%%% binaries, and allow for multiple variants to coexist, unique delimiters are
%%% used to distinguish RPM package names and module hierarchy.
%%% 
%%% In the case of a fresh install, \OHPC{} recipes default to installation of the
%%% latest toolchains available in a given release branch. However, if upgrading a
%%% previously installed system, administrators can {\em opt-in} to enable new
%%% variants as they become available. To illustrate this point, consider the
%%% previous \OHPC{} 1.3.5 release as an example which contained GCC 7.3.0
%%% along with runtimes and libraries compiled with this toolchain.  In the case
%%% where an admin would like to enable the newer {``gnu8''} toolchain,
%%% installation of these additions is simplified with the use of \OHPC{}'s
%%% meta-packages (see Table~\ref{table:groups} in Appendix
%%% \ref{appendix:manifest}).  The following example illustrates adding the
%%% complete ``gnu8'' toolchain.  Note that we leverage the convenience
%%% meta-packages containing MPI-dependent builds, and we also update the
%%% modules environment to make it the default.
%%% 
%%% \begin{lstlisting}[language=bash,keywords={}]
%%% # Install GCC 8.x-compiled meta-packages with dependencies
%%% [sms](*\#*)  (*\install*) ohpc-gnu8-perf-tools \
%%%                          ohpc-gnu8-io-libs \
%%%                          ohpc-gnu8-python-libs \
%%%                          ohpc-gnu8-runtimes \
%%%                          ohpc-gnu8-serial \
%%%                          ohpc-gnu8-parallel-libs
%%% 
%%% # Update default environment
%%% [sms](*\#*) (*\remove*) lmod-defaults-gnu7-openmpi3-ohpc
%%% [sms](*\#*) (*\install*) lmod-defaults-gnu8-openmpi3-ohpc
%%% 
%%% \end{lstlisting}
%%% 
%%% 

\input{common/test_suite}
\input{common/customization_appendix_centos}
\clearpage

\definecolor{Gray}{gray}{0.5}
\newcommand{\captionSpace}{-0.15cm}
\newcommand{\tabSpaceBot}{1.0cm}
\captionsetup{justification=raggedright,singlelinecheck=false}

\subsection{Package Manifest} \label {appendix:manifest}

\vspace*{0.25cm}
This appendix provides a summary of available meta-package groupings and all of
the individual RPM packages that are available as part of this \OHPC{}
release. The meta-packages provide a mechanism to group related collections of
RPMs by functionality and provide a convenience mechanism for installation.  A
list of the available meta-packages and a brief description is presented in
Table~\ref{table:groups}.

\vspace*{1.25cm}
\begin{table}[h!] 
\caption{\bf Available \OHPC{} Meta-packages} \vspace*{\captionSpace{}}
\label{table:groups}
\input manifest/patterns
\end{table}

% meta-packages (2)
\begin{table}[h!] 
\caption*{Table~\ref{table:groups} (cont): {\bf Available \OHPC{} Meta-packages} \vspace*{\captionSpace{}} }
\input manifest/patterns2
\end{table}

\iftoggleverb{isx86}
% meta-packages (3)
\begin{table}[h!] 
\caption*{Table~\ref{table:groups} (cont): {\bf Available \OHPC{} Meta-packages} \vspace*{\captionSpace{}} }
\input manifest/patterns3
\end{table}

\fi

\clearpage
What follows next in this Appendix is a series of tables that summarize the
underlying RPM packages available in this \OHPC{} release. These packages are
organized by groupings based on their general functionality and each table
provides information for the specific RPM name, version, brief summary, and the
web URL where additional information can be obtained for the component. Note
that many of the 3rd party community libraries that are pre-packaged
with \OHPC{} are built using multiple compiler and MPI families. In these cases,
the RPM package name includes delimiters identifying the development
environment for which each package build is targeted.  Additional information
on the \OHPC{} package naming scheme is presented in \S\ref{sec:3rdparty}. 
The relevant package groupings and associated Table references are as follows:

\vspace*{0.1cm}

\begin{itemize*}
\item Administrative tools (Table~\ref{table:admin})
\iftoggleverb{isWarewulf}
\item Provisioning (Table~\ref{table:provisioning})
\fi
\item Resource management (Table~\ref{table:rms})
\item Compiler families (Table~\ref{table:compiler-families})
\item MPI families (Table~\ref{table:mpi-families})
\item Development tools (Table~\ref{table:dev-tools})
\item Performance analysis tools (Table~\ref{table:perf-tools})
\iftoggleverb{isCentOS_x86}
\item Lustre (Table~\ref{table:lustre})
\fi

%\item Distro support packages and dependencies (Table~\ref{table:distro-packages})
\item IO Libraries (Table~\ref{table:io-libs})
\item Runtimes (Table~\ref{table:runtimes})
\item Serial/Threaded Libraries (Table~\ref{table:serial-libs})
\item Parallel Libraries (Table~\ref{table:parallel-libs})
\end{itemize*}


\newcommand{\firstColWidth}{5.2cm}
\newcommand{\secondColWidth}{1.25cm}

\vspace*{1.0cm}
\urlstyle{same}

% Administration Tools 
\begin{table}[h]
\caption{\bf Administrative Tools} \vspace*{\captionSpace{}} \label{table:admin}
\input manifest/admin
\end{table}
\vspace*{0.5cm}

\renewcommand{\firstColWidth}{6.25cm}
\renewcommand{\secondColWidth}{0.95cm}

% Provisioning

\begin{table}[h!]
\caption{\bf Provisioning} \vspace*{\captionSpace{}} \label{table:provisioning}
\input manifest/provisioning
\vspace*{\tabSpaceBot{}}
\end{table}

\renewcommand{\firstColWidth}{4.1cm}
\renewcommand{\secondColWidth}{1.8cm}

% Resource Management
\begin{table}[h!]
\caption{\bf Resource Management} \vspace*{\captionSpace{}} \label{table:rms}
\input manifest/rms
\vspace*{\tabSpaceBot{}}
\end{table}

\renewcommand{\firstColWidth}{4.7cm}

% Compiler Families
\begin{table}[h!]
\caption{\bf Compiler Families} \vspace*{\captionSpace{}} \label{table:compiler-families}
\input manifest/compiler-families
\vspace*{\tabSpaceBot{}}
\end{table}

\renewcommand{\firstColWidth}{4.7cm}

% MPI Families
\begin{table}[h!]
\caption{\bf MPI Families / Communication Libraries} \vspace*{\captionSpace{}} \label{table:mpi-families}
\input manifest/mpi-families
\vspace*{\tabSpaceBot{}}
\end{table}

\renewcommand{\firstColWidth}{5.6cm}
\renewcommand{\secondColWidth}{1.5cm}

% Development Tools
\begin{table}[h!]
\caption{\bf Development Tools} \vspace*{\captionSpace{}} \label{table:dev-tools}
\input manifest/dev-tools
\vspace*{\tabSpaceBot{}}
\end{table}
%%% 
%%% % Development Tools (2)
%%% \begin{table}[h!]
%%% \caption*{Table~\ref{table:dev-tools} (cont): {\bf Development Tools} \vspace*{\captionSpace{}} }
%%% \input manifest/dev-tools2
%%% \vspace*{\tabSpaceBot{}}
%%% \end{table}
%%% 
\renewcommand{\firstColWidth}{4.3cm}
\renewcommand{\secondColWidth}{3.1cm}

% Perf Tools (1)
\begin{table}[h!]
\caption{\bf Performance Analysis Tools} \vspace*{\captionSpace{}} \label{table:perf-tools}
\input manifest/perf-tools
\vspace*{\tabSpaceBot{}}
\end{table}

% Perf Tools (2)
\begin{table}[h!]
\caption*{Table~\ref{table:perf-tools} (cont): {\bf Performance Analysis Tools} \vspace*{\captionSpace{}} }
\input manifest/perf-tools2
\vspace*{\tabSpaceBot{}}
\end{table}

% Perf Tools (3)
\begin{table}[h!]
\caption*{Table~\ref{table:perf-tools} (cont): {\bf Performance Analysis Tools} \vspace*{\captionSpace{}} }
\input manifest/perf-tools3
\vspace*{\tabSpaceBot{}}
\end{table}

%%% % Perf Tools (4)
%%% \begin{table}[h!]
%%% \caption*{Table~\ref{table:perf-tools} (cont): {\bf Performance Analysis Tools} \vspace*{\captionSpace{}} }
%%% \input manifest/perf-tools4
%%% \vspace*{\tabSpaceBot{}}
%%% \end{table}
%%% 
%%% % Distro Packages
%%% \begin{table}[h!]
%%% \caption{\bf Distro Support Packages/Dependencies} \vspace*{\captionSpace{}} \label{table:distro-packages}
%%% \input manifest/distro-packages
%%% \vspace*{\tabSpaceBot{}}
%%% \end{table}
%%% 
\renewcommand{\firstColWidth}{5.1cm}
\renewcommand{\secondColWidth}{1.4cm}

%%% % Lustre
%%% \begin{table}[h!]
%%% \caption{\bf Lustre} \vspace*{\captionSpace{}} \label{table:lustre}
%%% \input manifest/lustre
%%% \vspace*{\tabSpaceBot{}}
%%% \end{table}
%%% 
%\renewcommand{\firstColWidth}{5.0cm}
%\renewcommand{\secondColWidth}{1.5cm}

% IO Libs (1)
\begin{table}[h!]
\caption{\bf IO Libraries} \vspace*{\captionSpace{}} \label{table:io-libs}
\input manifest/io-libs
\vspace*{\tabSpaceBot{}}
\end{table}

% IO Libs (2)
\begin{table}[h!]
\caption*{Table~\ref{table:io-libs} (cont): {\bf IO Libraries} \vspace*{\captionSpace{}} }
\input manifest/io-libs2
\vspace*{\tabSpaceBot{}}
\end{table}

%%% % IO Libs (3)
%%% \begin{table}[h!]
%%% \caption*{Table~\ref{table:io-libs} (cont): {\bf IO Libraries} \vspace*{\captionSpace{}} }
%%% \input manifest/io-libs3
%%% \vspace*{\tabSpaceBot{}}
%%% \end{table}
%%% 
% Runtimes
\clearpage
\begin{table}[h!]
\caption{\bf Runtimes} \vspace*{\captionSpace{}} \label{table:runtimes}
\input manifest/runtimes
\vspace*{\tabSpaceBot{}}
\end{table}

% Serial libs
\begin{table}[h!]
\caption{\bf Serial/Threaded Libraries} \vspace*{\captionSpace{}} \label{table:serial-libs}
\input manifest/serial-libs
\vspace*{\tabSpaceBot{}}
\end{table}

% Parallel libs (1)
\begin{table}[h!]
\caption{\bf Parallel Libraries} \vspace*{\captionSpace{}} \label{table:parallel-libs}
\input manifest/parallel-libs
\vspace*{\tabSpaceBot{}}
\end{table}

% Parallel libs (2)
\begin{table}[h!]
\caption*{Table~\ref{table:parallel-libs} (cont): {\bf Parallel Libraries} \vspace*{\captionSpace{}} }
\input manifest/parallel-libs2
\vspace*{\tabSpaceBot{}}
\end{table}

% Parallel libs (3)
\begin{table}[h!]
\caption*{Table~\ref{table:parallel-libs} (cont): {\bf Parallel Libraries} \vspace*{\captionSpace{}} }
\input manifest/parallel-libs3
\vspace*{\tabSpaceBot{}}
\end{table}

%%% % Parallel libs (4)
%%% \begin{table}[h!]
%%% \caption*{Table~\ref{table:parallel-libs} (cont): {\bf Parallel Libraries} \vspace*{\captionSpace{}} }
%%% \input manifest/parallel-libs4
%%% \vspace*{\tabSpaceBot{}}
%%% \end{table}
%%% 
%%% 

\input{common/signature}


\end{document}

